\section{An exact correction to a finite-field assertion}
\label{sec:audit}

The reviewed arXiv preprint of Radchenko--Zagier, version
\texttt{2012.15805v1}, Section~7.2, printed page~16, includes the
finite-field assertion
\begin{equation}\label{eq:incorrect-cocycle}
 \beta_p(n)=\beta_p(n+1)+\beta_p\bigl(n/(n+1)\bigr),
\end{equation}
where arguments are reduced modulo the prime $p$. The same statement
appears in the author-hosted preprint checked for this report
\cite{RZ}. For the exterior symbol defined in \eqref{eq:beta}, this
assertion is false. The issue is not a decimal discrepancy or an
unresolved branch convention: the obstruction is a nonzero rationalized
exterior symbol.

\begin{proposition}[Prime-seven counterexample]\label{prop:counterexample}
At $p=7$ and $n=2$, the defect in \eqref{eq:incorrect-cocycle} is
$3\beta_7(2)\ne0$.
\end{proposition}
\begin{proof}
Modulo $7$ one has $2/3=3$, and $3\equiv-2^{-1}$. The valid
symmetries \eqref{eq:symmetries} therefore give
$\beta_7(3)=-\beta_7(2)$. Equation
\eqref{eq:incorrect-cocycle} would say
$\beta_7(2)=-2\beta_7(2)$.

Here is an independent witness that $\beta_7(2)\ne0$, not relying on
the full kernel theorem. Write
\[
 a=\cl{1-\zeta_7},\quad b=\cl{1-\zeta_7^2},\quad
 c=\cl{1-\zeta_7^3}.
\]
Pairing $k$ and $-k$ in the definition gives
\begin{equation}\label{eq:beta7-wedge}
 \beta_7(2)=2(b\wedge a+c\wedge b+a\wedge c)
          =-2(b-a)\wedge(c-a).
\end{equation}
Let
\[
 A=\log(2\sin(\pi/7)),\quad B=\log(2\sin(2\pi/7)),\quad
 C=\log(2\sin(3\pi/7)).
\]
Since $0<\pi/7<2\pi/7<3\pi/7<\pi/2$, one has $A<B<C$.
Take logarithmic absolute values at the two embeddings indexed by
$1$ and $2$. The vectors corresponding to $b-a$ and $c-a$ are
respectively
\[
 (B-A,C-B),\qquad(C-A,A-B).
\]
Their determinant is
\begin{equation}\label{eq:det-witness}
 -(B-A)^2-(C-B)(C-A)<0.
\end{equation}
Thus their exterior product has a nonzero image, proving
\eqref{eq:beta7-wedge} nonzero. The defect is consequently
$3\beta_7(2)$, as claimed.
\end{proof}

\subsection{Precise correction and unaffected statements}

For the definition in \eqref{eq:beta}, the universal finite-field
three-term identity \eqref{eq:incorrect-cocycle} should be removed.
Any intended replacement involving an altered module, action or quotient
would need a new statement and proof. It is not justified by applying
rational Herglotz three-term relations directly to finite-field residues.

The sign and inverse symmetries remain true. The rational boundary
formula \eqref{eq:boundary-F} also remains true: it follows directly
from finite products, as proved in Proposition~\ref{prop:boundary}.
The current ProveIt manuscript already includes such a direct derivation
and explicitly separates it from a finite-field cocycle claim. The
new classification therefore does not inherit the erroneous premise.

This correction is deliberately version-specific. It identifies the
statement in the preprint pages actually inspected; it does not claim
that all editions of the journal article, all subsequent corrections,
or every cocycle construction in the paper have been checked. The
analytic Hecke relation and the finite dilogarithm evaluations used in
Section~\ref{sec:identities} do not use
\eqref{eq:incorrect-cocycle} in their displayed proofs.

\subsection{Editorial changes recommended for the manuscript}

The next integration can replace the open general-$J(p/q)$ formal
classification question with Theorem~\ref{thm:classification}, while
retaining the separate numerical period question. The phrase
``modulo $2e$'' should be preserved in every summary of the mixed-parity
criterion; the pair $(8,3)$ is a useful regression test.

The analytic identity \eqref{eq:neighbor} and the evaluation
\eqref{eq:J35} can be placed immediately after the existing $J(2/q)$
section. The finite-field caveat belongs beside the direct boundary
calculation. The theorem ledger should mark the existing kernel theorem,
the $K$-theory facts and the Hecke relation as prior inputs, not new
results. No earlier Gaussian or $S_p$ conjecture should have its status
changed as a side effect of this integration.
